% !TEX root = ../notes.tex

\documentclass[../notes.tex]{subfiles}

\begin{document}

\section{April 8}
Here we go.

\subsection{The Deformation Theorem}
We want to find a universal deformation. To begin, we should give notation for our moduli space.
\begin{notation}
	For a connective $\mathbb E_\infty$-ring $R$, we let $\mc M_{\mathrm{BT}}(R)$ denote the groupoid of $p$-divisible groups on $R$.
\end{notation}
Here, $\mathrm{BT}$ stands for Barsotti--Tate. It is not clear what Barsotti did to earn a name here.
\begin{notation}
	Choose a $p$-divisible group $\mathbb G_0$ over a discrete ring $R_0$. Then let $X$ be the functor from connective $\mathbb E_\infty$-rings to spaces defined by
	\[X(A)\coloneqq\colim_{\substack{I\subseteq\pi_0A\\I\text{ nilpotent}}}\mc M_{\mathrm{BT}}(A)\times_{\mc B_{\mathrm{BT}}(\pi_0(A/I))}\op{Hom}_{\mathrm{Ring}}(R_0,\pi_0(A/I)).\]
	Here, $\mathbb G_0$ appears in the fiber product in the right-hand term.
\end{notation}
\begin{remark}
	Intuitively, $X$ classifies deformations of the pair $(R_0,\mathbb G_0)$ to $A$. The moral is that we are hunting for $p$-divisible groups $\mathbb G$ over $A$ which suitably descend to $\mathbb G_0$ under a given map $R_0\to\pi_0(A/I)$.
\end{remark}
We are now ready to state our target theorem.
\begin{theorem}[Lurie] \label{thm:lurie}
	Choose a $p$-divisible group $\mathbb G_0$ over a discrete ring $R_0$.
	Suppose that $p$ is nilpotent in $R_0$, $L_{R_0}$ is an almost perfect $R_0$-module, and $\mathbb G_0$ is non-stationary.
	Then there is a connective $\mathbb E_\infty$-ring $R_0^{\mathrm{def}}$ such that $X\coloneqq\op{Spf}R_0^{\mathrm{def}}$, where we are taking $\op{Spf}$ along the kernel of $\pi_0R_0^{\mathrm{def}}\to R_0$.
\end{theorem}
\begin{remark}
	If $R_0$ is a perfect algebra over $\FF_p$, then the conditions are all automatic. Here, an $\mathbb F_p$-algebra is perfect if and only if the Frobenius is bijective.
\end{remark}
\begin{remark}
	More directly, a map $R_0^{\mathrm{def}}\to A$ of $\mathbb E_\infty$-rings parameterizes $p$-divisible groups $\mathbb G$ over $A$, together with a map $f\colon A\to R_0$ and an isomorphism $f_*\mathbb G\cong\mathbb G_0$.
\end{remark}
\begin{remark}
	It turns out that $\pi_0R_0^{\mathrm{def}}$ has a universal property in classical commutative algebra. The existence of this classical object is the hardest part of the proof.
\end{remark}
\begin{example}
	If $R_0=\FF_p$ and $\mathbb G_0=\widehat{\mathbb G}_m$, then it will turn out that $R_0^{\mathrm{def}}=\mathbb S_p^\land$.
\end{example}
We are owed some definitions. Let's start with $L_{R_0}$. Intuitively, $L_{R_0}$ keeps tracks of deformations of the ring $R_0$, so we should start by finding such a deformation.
\begin{definition}
	Fix an $\mathbb E_\infty$-ring $R$, and let $M$ be an $R$-module. Then the \textit{trivial square-zero extension} is an $\mathbb E_\infty$-algebra $R\oplus M$. Multiplication is given by $(_1,m_1)\cdot(r_2,m_2)=(r_1r_2,r_1m_2+r_2m_1)$.
\end{definition}
\begin{remark}
	This extension has the ideal $M\subseteq R\oplus M$ being square-zero. Also, $R\oplus M$ is an extension of $R$, where the kernel is this ideal.
\end{remark}
\begin{remark}
	Intuitively, $R\oplus M$ is the most boring possible extension of $R$.
\end{remark}
\begin{remark}
	In general, a square-zero extension of $R$ is given by the pullback of the diagram
	% https://q.uiver.app/#q=WzAsNSxbMSwwLCJSIl0sWzAsMSwiUiJdLFsxLDEsIlJcXG9wbHVzIE0iXSxbMiwwLCJyIl0sWzIsMSwiKHIsMCkiXSxbMSwyXSxbMCwyXSxbMyw0LCIiLDIseyJzdHlsZSI6eyJ0YWlsIjp7Im5hbWUiOiJtYXBzIHRvIn19fV1d&macro_url=https%3A%2F%2Fraw.githubusercontent.com%2FdFoiler%2Fnotes%2Fmaster%2Fnir.tex
	\[\begin{tikzcd}[cramped]
		& R & r \\
		R & {R\oplus M} & {(r,0)}
		\arrow[from=1-2, to=2-2]
		\arrow[maps to, from=1-3, to=2-3]
		\arrow[from=2-1, to=2-2]
	\end{tikzcd}\]
	where the bottom map is allowed to be arbitrary. Notably, in the case where the bottom map is $r\mapsto(r,0)$, we get the trivial square-zero extension, but there may be more: for example, there is a fiber sequence $\ZZ/p^2\ZZ\to\FF_p\to\Sigma\FF_p$, so $\ZZ/p^2\ZZ$ is a possible square-zero extension of $\FF_p$.
\end{remark}
The previous remark motivates the following definition.
\begin{definition}[cotangent complex]
	Fix an $\mathbb E_\infty$-ring $R$. Roughly speaking, we define the left adjoint of the trivial square zero extension (from $R$-modules to $\mathbb E_\infty$-algebras) to be $L_R$. Explicitly,
	\[\op{Hom}_{\op{Alg}_R(\mathbb E_\infty)}(R,R\oplus M)=\op{Hom}_{\mathrm{Mod}(R)}(L_R,M).\]
	Note that $L_R$ is not technically a left adjoint here because the ambient categories themselves depend on $R$.
\end{definition}
\begin{remark}
	Intuitively, an $M$-point of $L_R$ (namely, a map $L_R\to M$) should look like a space of square-zero maps $R\to M$, which is roughly speaking a cotangent vector.
\end{remark}
\begin{remark}
	Note that square-zero extensions are classified by the ring maps $R\to R\oplus M$, which we know see is parameterized by maps $L_R\to M$. Thus, $L_R$ knows how to deform the ring $R$!
\end{remark}
\begin{example}
	Suppose $R$ is an $\mathbb E_\infty$-ring and $n\ge0$. Then $\tau_{\le n+1}R$ is a square-zero extension of $\tau_{\le n}R$. Indeed, there is a fiber sequence
	\[\Sigma^{n+1}\pi_{n+1}R\to\tau_{\le n+1}R\to\tau_{\le n}R\to\Sigma^{n+2}\pi_{n+1}R\]
	of spectra, which induces a pullback square
	% https://q.uiver.app/#q=WzAsNCxbMCwwLCJcXHRhdV97XFxsZSBuKzF9UiJdLFsxLDAsIlxcdGF1X3tcXGxlIG59UiJdLFsxLDEsIlxcdGF1X3tcXGxlIG59Ulxcb3BsdXNcXFNpZ21hXntuKzF9XFxwaV97bisxfVIiXSxbMCwxLCJcXHRhdV97XFxsZSBufVIiXSxbMCwxXSxbMSwyXSxbMywyXSxbMCwzXV0=&macro_url=https%3A%2F%2Fraw.githubusercontent.com%2FdFoiler%2Fnotes%2Fmaster%2Fnir.tex
	\[\begin{tikzcd}[cramped]
		{\tau_{\le n+1}R} & {\tau_{\le n}R} \\
		{\tau_{\le n}R} & {\tau_{\le n}R\oplus\Sigma^{n+1}\pi_{n+1}R}
		\arrow[from=1-1, to=1-2]
		\arrow[from=1-1, to=2-1]
		\arrow[from=1-2, to=2-2]
		\arrow[from=2-1, to=2-2]
	\end{tikzcd}\]
	in spectra, where the bottom map uses the induced map $\tau_{\le n}R\to\Sigma^{n+2}\pi_{n+1}R$, but the right map only includes into the left factor; in fact, these morphisms upgrade to ring morphisms, so this is a pullback in $\mathbb E_\infty$-rings.
\end{example}
Now, for \Cref{thm:lurie}, we would like for $L_{R_0}$ to be almost perfect.
\begin{definition}[almost perfect]
	Fix a connective module $M$ over an $\mathbb E_\infty$-ring $R$. Then $M$ is \textit{perfect} if and only if it is compact, meaning that $\op{Hom}_R(M,-)$ commutes with filtered colimits. Similarly, $M$ is \textit{almost perfect} if and only if $\op{Hom}_R(M,-)$ commutes with filtered colimits over diagrams $\{N_i\}$ for which the $N_i$s have uniformly bounded homotopy groups.
\end{definition}
\begin{remark}
	If $R$ is a discrete ring so that $\mathrm{Mod}(R)$ can be identified with chain complexes, then being perfect means being in the thick subcategory generated by $R$. This turns out to be equivalent to being compact in $\mathrm{Mod}(R)$; for example, certainly any power of $R$ is compact.
\end{remark}
Lastly, we are owed a definition of non-stationary.
\begin{definition}[non-stationary]
	Fix a $p$-divisible group $\mathbb G_0$ over a discrete ring $R_0$. Then $\mathbb G_0$ is \textit{non-stationary} if and only if the following holds: for every point $x\in\Spec R_0$ and derivation $d\colon R_0\to k(x)$, there is an associated square-zero extension
	\[R_0\to R_0\oplus k(x)\to R_0\]
	whose associated $p$-divisible group is a nontrivial deformation of $\mathbb G_0$.
\end{definition}
\begin{remark}
	Intuitively, we are asking for $\mathbb G_0$ to deform when $R_0$ deforms.
\end{remark}
Let's explain the relevance of these notions for \Cref{thm:lurie}. Note that the $p$-divisible group $\mathbb G_0$ over $R_0$ amounts to a map $\Spec R_0\to\mc M_{\mathrm{BT}}$. Then we define a relative cotangent complex $L_{R_0/\mc M_{\mathrm{BT}}}$, and it will be almost perfect if and only if $L_{R_0}$ is almost perfect. Then we will show that $\pi_*L_{R_0/\mc M_{\mathrm{BT}}}$ is concentrated in degrees at least one if and only if $\mathbb G_0$ is non-stationary. From here, these conditions combine with Artin representability to receive a formal scheme.

We are now motivated to define a relative cotangent complex.
\begin{definition}[quasicoherent sheaves]
	Fix a functor $X$ from connective $\mathbb E_\infty$-rings to spaces. Then the category of \textit{quasicoherent sheaves} is
	\[\op{QCoh}X\coloneqq\lim_{\Spec A\to X}\mathrm{Mod}(A).\]
\end{definition}
\begin{definition}[relative cotangent complex]
	Fix functors $X$ and $Y$ from connective $\mathbb E_\infty$-rings to the category of spaces, and choose a natural transformation $f\colon X\to Y$. Then we define $L_{X/Y}$ (provided it exists) to be a quasicoherent sheaf on $X$ such that, for any $\eta\colon\Spec A\to X$ and map $\Spec(A\oplus M)\to Y$, we have that $\op{Hom}_A(\eta^*L_{X/Y},M)$ classifies dashed arrows fitting into the following diagram.
	% https://q.uiver.app/#q=WzAsNCxbMCwwLCJcXFNwZWMgQSJdLFsxLDAsIlgiXSxbMCwxLCJcXFNwZWMoQVxcb3BsdXMgTSkiXSxbMSwxLCJZIl0sWzAsMSwiXFxldGEiXSxbMSwzLCJmIl0sWzIsM10sWzAsMl0sWzIsMSwiIiwxLHsic3R5bGUiOnsiYm9keSI6eyJuYW1lIjoiZGFzaGVkIn19fV1d&macro_url=https%3A%2F%2Fraw.githubusercontent.com%2FdFoiler%2Fnotes%2Fmaster%2Fnir.tex
	\[\begin{tikzcd}[cramped]
		{\Spec A} & X \\
		{\Spec(A\oplus M)} & Y
		\arrow["\eta", from=1-1, to=1-2]
		\arrow[from=1-1, to=2-1]
		\arrow["f", from=1-2, to=2-2]
		\arrow[dashed, from=2-1, to=1-2]
		\arrow[from=2-1, to=2-2]
	\end{tikzcd}\]
\end{definition}

\end{document}